%&pdflatex
\documentclass[a4paper,12pt]{article}

\addtolength{\textwidth}{1cm}

\usepackage[english,spanish]{babel}

\usepackage[leqno]{amsmath}
\usepackage{amsfonts,amsthm}

\nofiles %% .aux files unnecessary

\theoremstyle{plain}
\newtheorem*{theoremOne}{Theorem 1}
\newtheorem*{propTwo}{Proposition 2}
\newtheorem*{corThree}{Corollary 3}
\newtheorem*{corFour}{Corollary 4}
\newtheorem*{theoremFive}{Theorem 5}

\theoremstyle{definition}
\newtheorem*{note}{Note}

\newcommand{\N}{\mathbb{N}}

\date{}

%%%%%%%%%%%%%%%%%%%%%%%%%%%%%%%%%%%%%%%%%%%%%%%%%%%%%%%%%%%%%%%
\title{\large \bfseries \boldmath $L^{p}$-BOUNDEDNESS OF THE
KERNELS RELATIVE TO GENERALIZED JACOBI WEIGHTS}

\author{%
J. J. Guadalupe Hern\'{a}ndez$^{(1)}$, M. P\'{e}rez Riera$^{(1)}$,\\[2pt]
F. J. Ruiz Blasco$^{(1)}$ and J. L. Varona Malumbres$^{(2)}$\\[5pt]
\small $^{(1)}$ Departamento de Matem\'{a}ticas,\\[-2pt]
\small Facultad de Ciencias, Universidad de Zaragoza.\\[3pt]
\small $^{(2)}$ Departamento de Matem\'{a}tica Aplicada,\\[-2pt]
\small Colegio Universitario de La Rioja, Universidad de Zaragoza.%
\renewcommand\thefootnote{}%
\renewcommand\footnotemark{}%
\footnote{This paper were published in
``Polinomios Ortogonales y Aplicaciones''
(Actas del VI Simposium, Gij\'on, 1989),
168--178, Universidad de Oviedo, Gij\'on, 1990.}%
}
%%%%%%%%%%%%%%%%%%%%%%%%%%%%%%%%%%%%%%%%%%%%%%%%%%%%%%%%%%%%%%%

\begin{document}

%%%%%%%%%%%%%%%%%%%%%%%%%%%%%%%%%%%%%%%%%%%%
\maketitle
%%%%%%%%%%%%%%%%%%%%%%%%%%%%%%%%%%%%%%%%%%%%

\noindent\textit{A.M.S. Classification}: 42C10.

\selectlanguage{spanish}

\begin{abstract}
Sea $w$ un peso de Jacobi generalizado sobre el intervalo $[-1,1]$,
es decir, $w(x) = h(x)(1-x)^{\alpha} (1+x)^{\beta}
\prod_{i=1}^{N} |x-x_{i}|^{\gamma_{i}}$, con $\alpha, \beta,
\gamma_{i} > -1$ y ciertas condiciones de continuidad
sobre~$h$. Mediante el uso de la teor\'{\i}a de pesos $A_{p}$, se puede
demostrar la convergencia de la serie de Fourier de polinomios
ortonormales con respecto a $w$ cuando $\gamma_{i} \geq 0$
$\forall i$. En este trabajo obtenemos acotaciones
de las normas en $L^{p}$ de los n\'{u}cleos relativos a $w$, que permiten
extender el resultado de la convergencia en media al caso general.
\end{abstract}

\selectlanguage{english}

%%%%%%%%%%%%%%%%%%%%%%%%%%%%%%%%%%%%%%%%%%%%

Let $w$ be a generalized Jacobi weight, that is,
$$
w(x) = h(x)(1-x)^{\alpha} (1+x)^{\beta}
\prod_{i=1}^{N} |x-x_{i}|^{\gamma_{i}},
\qquad x \in [-1,1],
$$
where:
\begin{itemize}
\item[(a)] $\alpha, \beta, \gamma_{i} > -1$,
$t_{i} \in (-1,1)$, $t_{i}\neq t_{j}$ $\forall i \neq j$;
\item[(b)] $h$ is a positive, continuous function on $[-1,1]$ and
$\omega(h,\delta) \delta^{-1} \in L^{1}(0,2)$,
$\omega(h,\delta)$ being the modulus of continuity of~$h$.
\end{itemize}

Let $d\mu = w(x)\,dx$ on $[-1,1]$ and let $S_{n}$ ($n \geq 0$)
be the $n$-th partial sum of the Fourier series in the orthonormal
polynomials with respect to $d\mu$. The study of the boundedness
\begin{equation}
\| S_{n}f \|_{L^{p}(u^{p}\,d\mu)}
\leq C \| f \|_{L^{p}(v^{p}\,d\mu)},
\tag{1}
\end{equation}
where
$$
u(x) = (1- x)^{a}(1+ x)^{b} \prod_{i=1}^{N} | x- t_{i} |^{g_{i}}
$$
and
$$
v(x) = (1- x)^{A}(1+ x)^{B} \prod_{i=1}^{N} | x- t_{i} |^{G_{i}}
$$
was done by Badkov ([B]) in the case $u = v$
by means of a direct estimation of the kernels $K_{n}(x,y)$
associated with the polynomials orthogonal with respect to $d\mu$.
Later, Varona ([V]) considered the same problem, with $u$ and $v$ not
necessarily equal; his method consists of an appropriate use
of the theory of $A_{p}$ weights. He found conditions for (1)
which generalized those obtained for $u = v$
by Badkov. However, this result, which we state below, follows
only in the case $\gamma_{i} \geq 0$ $\forall i$:

\begin{theoremOne}
Let $\gamma_{i} \geq 0$ $\forall i$ and $1 < p < \infty$.
If the inequalities
\begin{equation}
\left\{
\begin{aligned}
A + (\alpha+1) \left(\frac{1}{p} - \frac{1}{2}\right) &< \frac{1}{4}, \\
B + (\beta+1) \left(\frac{1}{p} - \frac{1}{2}\right) &< \frac{1}{4}, \\
G_{i} + (\gamma_{i}+1) \left(\frac{1}{p} - \frac{1}{2}\right) &< \frac{1}{2}
\quad \forall i,
\end{aligned}
\right.
\tag{2}
\end{equation}
\begin{equation}
\left\{
\begin{aligned}
A + (\alpha+1) \left(\frac{1}{p} - \frac{1}{2}\right) &< \frac{\alpha+1}{2}, \\
B + (\beta+1) \left(\frac{1}{p} - \frac{1}{2}\right) &< \frac{\beta+1}{2}, \\
G_{i} + (\gamma_{i}+1) \left(\frac{1}{p} - \frac{1}{2}\right) &< \frac{\gamma_{i}+1}{2}
\quad \forall i,
\end{aligned}
\right.
\tag{3}
\end{equation}
\begin{equation}
\left\{
\begin{aligned}
a + (\alpha+1) \left(\frac{1}{p} - \frac{1}{2}\right) &> - \frac{1}{4}, \\
b + (\beta+1) \left(\frac{1}{p} - \frac{1}{2}\right) &> - \frac{1}{4}, \\
g_{i} + (\gamma_{i}+1) \left(\frac{1}{p} - \frac{1}{2}\right) &> - \frac{1}{2}
\quad \forall i,
\end{aligned}
\right.
\tag{4}
\end{equation}
\begin{equation}
\left\{
\begin{aligned}
a + (\alpha+1) \left(\frac{1}{p} - \frac{1}{2}\right) &> - \frac{\alpha+1}{2}, \\
b + (\beta+1) \left(\frac{1}{p} - \frac{1}{2}\right) &> - \frac{\beta+1}{2}, \\
g_{i} + (\gamma_{i}+1) \left(\frac{1}{p} - \frac{1}{2}\right) &> - \frac{\gamma_{i}+1}{2}
\quad \forall i,
\end{aligned}
\right.
\tag{5}
\end{equation}
\begin{equation}
A \leq a, \qquad B \leq b, \qquad G_{i} \leq g_{i}
\quad \forall i,
\tag{6}
\end{equation}
hold, then there exists $C > 0$ such that
$$
\| S_{n}f \|_{L^{p}(u^{p}\,d\mu)}
\leq C \| f \|_{L^{p}(v^{p}\,d\mu)}
\quad \forall f \in L^{p}(v^{p}\,d\mu),
\quad \forall n \in \N.
$$
\end{theoremOne}

The objective of this paper is to eliminate the restriction $\gamma_{i} \geq 0$,
by studying the norms of the kernels $K_{n}(x,y)$. We are going
to use the following notation: $\{P_{n}(x)\}$, $\{k_{n}\}$ and
$\{K_{n}(x,y)\}$ will be, respectively, the orthonormal polynomials,
their leading coefficients and the kernels relatives to $d\mu$;
if $c \in (-1,1)$, $\{P_{n}^{c}(x)\}$,
$\{k_{n}^{c}\}$ and $\{K_{n}^{c}(x,y)\}$ will be the corresponding
to $(x-c)^{2}\,d\mu$. Then, it is not difficult to establish
the following relations, $\forall n \geq 1$:
\begin{equation}
K_{n}(x,y) = (x-c) (y-c) K_{n-1}^{c}(x,y)
+ \frac{K_{n}(x,c)K_{n}(c,y)}{K_{n}(c,c)}\,,
\tag{7}
\end{equation}
\begin{equation}
K_{n}(x,c) = \frac{k_{n}}{k_{n}^{c}}\, P_{n}(c) P_{n}^{c}(x)
- \frac{k_{n-1}^{c}}{k_{n+1}}\, P_{n+1}(c) P_{n-1}^{c}(x).
\tag{8}
\end{equation}
It can be also shown (see [MNT], theorems~10 and~11, and~[R],
page~212) that
\begin{equation}
\lim_{n \to \infty} \frac{k_{n}}{k_{n}^{c}}
= \lim_{n \to \infty} \frac{k_{n-1}^{c}}{k_{n+1}}
= \frac{1}{2}.
\tag{9}
\end{equation}

If we define
$$
d(x,n) = (1 - x + n^{-2})^{-(2\alpha+1)/4}
(1 + x + n^{-2})^{-(2\beta+1)/4}
\prod_{i=1}^{N} ( | x-t_{i}| + n^{-1})^{-\gamma_{i}/2},
$$
it is known ([B]) that there exists a constant $C$ such that
$\forall x \in [-1,1]$, $\forall n$,
\begin{equation}
| P_{n}(x) | \leq C d(x,n).
\tag{10}
\end{equation}

There are also some well-known estimates for the kernels, one
of them being this ([N], page~4 and page~119, theorem~25): if
$c \in (-1,1)$ and the factor $| x - c |$
is present in $w$ with an exponent $\gamma$, there exist some
positive constants $C_{1}$ and $C_{2}$, depending on $c$,
such that,~$\forall n$,
\begin{equation}
C_{1} n^{\gamma+1} \leq K_{n}(c,c) \leq C_{2} n^{\gamma+1}.
\tag{11}
\end{equation}

From now on, all constants will be denoted $C$, so by $C$ we will
mean a constant, possibly different in each occurrence. Using
(8), (9) and (10), we get the following result:

\begin{propTwo}
Let $1 < p < \infty$, $1/p + 1/q = 1$.
Suppose the inequalities (4) and (5) hold. Let $-1 < c < 1$
and let $\gamma$ and $g$ be the exponents of $| x - c |$
in $w$ and $u$, respectively. Then, there exists a positive constant
$C$ such that, $\forall n \geq 0$,
$$
\| K_{n}(x,c) \|_{L^{p}(u^{p}w)}
\leq
\begin{cases}
C n^{(\gamma+1)/q - g}, &\text{if}\quad g < (\gamma +1) (1/2 - 1/p) + 1/2,\\
C n^{\gamma/2} (\log n)^{1/p}, &\text{if}\quad g = (\gamma +1) (1/2 - 1/p) + 1/2,\\
C n^{\gamma/2}, &\text{if}\quad (\gamma +1) (1/2 - 1/p) + 1/2 < g.
\end{cases}
$$
\end{propTwo}

\begin{proof}
From (10) it follows that $| P_{n}(c) | \leq C n^{\gamma/2}$.
Since $\{P_{n}^{c}\}$ is the sequence
associated with $(x-c)^{2}\,d\mu$, it also follows from (10)
that
$$
| P_{n}^{c}(x) | \leq C ( | x - c | + n^{-1})^{-1} d(x,n).
$$
Now, from (8) and (9) we get
\begin{equation}
| K_{n}(x,c) | \leq C n^{\gamma/2} ( | x - c | + n^{-1})^{-1} d(x,n).
\tag{12}
\end{equation}
Let us take $\varepsilon > 0$ such that $| t_{i} - c | > \varepsilon$
$\forall i$ for all $t_{i} \neq c$.
We can write
\begin{multline*}
\| K_{n}(x,c) \|^{p}_{L^{p}(u^{p}w)} \\
= \int_{| x - c | > \varepsilon} | K_{n}(x,c) |^{p} u(x)^{p} w(x)\,dx
+ \int_{| x - c | < \varepsilon} | K_{n}(x,c) |^{p} u(x)^{p} w(x)\,dx.
\end{multline*}
Using (12), for the first term we obtain
\begin{align*}
&\int_{| x - c | > \varepsilon} | K_{n}(x,c) |^{p} u(x)^{p} w(x)\,dx \\
&\qquad\leq C n^{p \gamma/2} \int_{| x - c | > \varepsilon}
( | x - c | + n^{-1})^{-p} d(x,n)^{p} u(x)^{p} w(x)\,dx \\
&\qquad\leq C_{1} n^{p \gamma/2} \int_{-1}^{1} d(x,n)^{p} u(x)^{p} w(x)\,dx.
\end{align*}
It is easy to deduce from (4) and (5) that this last integral
is bounded by a constant which does not depend on $n$, so
\begin{equation}
\int_{| x - c | > \varepsilon} | K_{n}(x,c) |^{p} u(x)^{p} w(x)\,dx
\leq C n^{p\gamma/2}.
\tag{13}
\end{equation}
Let us take now the second term; since for $| x - c | < \varepsilon$
there exists a constant $C$ such that, $\forall n$,
$d(x,n) \leq C ( | x - c | + n^{-1})^{-\gamma/2}$,
$u(x) \leq C | x - c |^{g}$
and $w(x) \leq C | x - c |^{\gamma}$, we have
\begin{align*}
&\int_{| x - c | < \varepsilon} | K_{n}(x,c) |^{p} u(x)^{p} w(x)\,dx \\
&\qquad\leq C n^{p\gamma/2} \int_{| x - c | < \varepsilon} ( | x - c | + n^{-1})^{-p}
d(x,n)^{p} u(x)^{p} w(x)\,dx \\
&\qquad\leq C_{1} n^{p \gamma/2}
\int_{| x - c | < \varepsilon} ( | x - c | + n^{-1})^{-p(1+\gamma/2)}
| x - c |^{gp+\gamma} \,dx \\
&\qquad\leq C _{2} n^{p\gamma/2}
\int_{0}^{1} (y + n^{-1})^{-p(1+\gamma/2)} y^{gp+\gamma} \,dy \\
&\qquad= C_{2} n^{p\gamma/2+p(1+\gamma/2)-gp-\gamma-1}
\int_{0}^{1} (ny + 1)^{-p(1+\gamma/2)} (ny)^{gp+\gamma} n\,dy \\
&\qquad= C_{2} n^{p\gamma/2+p(1+\gamma/2)-gp-\gamma-1}
\int_{0}^{n} (r + 1)^{-p(1+\gamma/2)} r^{gp+\gamma} \,dr.
\end{align*}
Taking into account that
$p (1+\gamma/2) - gp - \gamma - 1 = p [ (\gamma+1) (1/2 - 1/p) - g + 1/2]$
and there exist some constants $C_{1}$ and $C_{2}$ such that
$C_{1} \leq r+1 \leq C_{2}$ on $[0,1]$ and $C_{1}r \leq r+1 \leq C_{2}r$
on $[1,n]$, we finally get the inequality
\begin{equation}
\begin{aligned}
&\int_{| x - c | < \varepsilon} | K_{n}(x,c) |^{p} u(x)^{p} w(x)\,dx \\
&\qquad\leq C n^{p \gamma/2+p [ (\gamma+1)(1/2 - 1/p) - g+1/2 ]}
\int_{0}^{1} r^{gp+\gamma} \,dr \\
&\qquad\qquad\null+ C n^{p\gamma/2+p [ (\gamma+1)(1/2 - 1/p) - g+1/2 ]}
\int_{1}^{n} r^{-p [ (\gamma+1)(1/2 - 1/p) - g+1/2 ] -1} \,dr.
\end{aligned}
\tag{14}
\end{equation}
Since (5) implies $gp + \gamma > -1$,
the first term can be bounded in this way:
\begin{equation}
C n^{p\gamma/2+p [ (\gamma+1)(1/2 - 1/p) - g+1/2 ]}
\int_{0}^{1} r^{gp+\gamma} \,dr
\leq C n^{p \gamma/2+p [ (\gamma+1)(1/2 - 1/p) - g+1/2 ]}.
\tag{15}
\end{equation}
Let us consider separately the three cases in the statement:

    (a) If $g < (\gamma +1) (1/2 - 1/p) + 1/2$,
then $-p [ (\gamma +1) (1/2 - 1/p) -g+ 1/2 ] -1 < -1$; so:
$$
\int_{1}^{n} r^{-p [ (\gamma+1)(1/2 - 1/p) - g+1/2 ] -1} \,dr \leq C.
$$
In this case, (14) and (15) imply
$$
\int_{| x - c | < \varepsilon} | K_{n}(x,c) |^{p} u(x)^{p} w(x)\,dx
\leq C n^{p\gamma/2+p [ (\gamma+1)(1/2 - 1/p) - g+1/2 ]}.
$$
Since $p [ (\gamma +1)(1/2 - 1/p) - g+1/2 ] > 0$,
from this inequality and (13) we obtain
\begin{align*}
\| K_{n}(x,c) \|^{p}_{L^{p}(u^{p}w)}
&\leq C n^{p\gamma/2+p [ (\gamma+1)(1/2 - 1/p) - g+1/2 ]} \\
&= C n^{p [ (\gamma+1)(1 - 1/p) - g ]}
= C n^{p [ (\gamma+1)/q - g ]},
\end{align*}
as we had to prove.

    (b) If $(\gamma +1) (1/2 - 1/p) + 1/2 < g$,
then $-p [ (\gamma +1) (1/2 - 1/p) -g+ 1/2 ] -1 > -1$; therefore:
$$
\int_{1}^{n} r^{-p [ (\gamma+1)(1/2 - 1/p) - g+1/2 ] -1} \,dr
\leq C n^{-p [ (\gamma+1)(1/2 - 1/p) - g+1/2 ] }.
$$
From (14) and (15), we get now
$\int_{| x - c | < \varepsilon} | K_{n}(x,c) |^{p} u(x)^{p} w(x)\,dx
\leq C n^{p\gamma/2}$
and, finally,
$$
\| K_{n}(x,c) \|^{p}_{L^{p}(u^{p}w)} \leq C n^{p\gamma/2}.
$$

    (c) If $g = (\gamma +1) (1/2 - 1/p) + 1/2$,
$$
\int_{1}^{n} r^{-p [ (\gamma+1)(1/2 - 1/p) - g+1/2 ] -1} \,dr \leq C \log n;
$$
hence,
$$
\int_{| x - c | < \varepsilon} | K_{n}(x,c) |^{p} u(x)^{p} w(x)\,dx
\leq C n^{p\gamma/2} \log n
$$
and
$$
\| K_{n}(x,c) \|^{p}_{L^{p}(u^{p}w)} \leq C n^{p\gamma/2} \log n.
$$
This concludes the proof of the proposition.
\end{proof}

\begin{corThree}
Let $1 < p < \infty$, $1/p + 1/q = 1$.
Suppose the inequalities (2) and (3) hold. Let $-1 < c < 1$
and let $\gamma$ and $G$ be the exponents of $| x - c |$
in $w$ and $v$, respectively. Then, there exists a positive constant
$C$ such that, $\forall n \geq 0$,
$$
\| K_{n}(x,c) \|_{L^{q}(v^{-q}w)}
\leq
\begin{cases}
C n^{\gamma/2}, &\text{if}\quad G < (\gamma +1) (1/2 - 1/p) - 1/2,\\
C n^{\gamma/2} (\log n)^{1/q}, &\text{if}\quad G = (\gamma +1) (1/2 - 1/p) - 1/2,\\
C n^{(\gamma+1)/p + G}, &\text{if}\quad (\gamma +1) (1/2 - 1/p) - 1/2 < G.
\end{cases}
$$
\end{corThree}

\begin{proof}
Just apply proposition~2 to the weight $v^{-1}$ and keep in mind
the equality $1/2 - 1/p = 1/q - 1/2$.
\end{proof}

The following result is just what we need to extend theorem~1
to the general case $\gamma_{i} > -1$:

\begin{corFour}
Let $1 < p < \infty$, $1/p + 1/q = 1$.
Suppose the inequalities (2), (3), (4), (5) and (6) hold. Let $-1 < c < 1$.
Then, there exists a positive constant $C$ such that, $\forall n \geq 0$,
$$
\| K_{n}(x,c) \|_{L^{p}(u^{p}w)} \| K_{n}(x,c) \|_{L^{q}(v^{-q}w)}
\leq C K_{n}(c,c).
$$
\end{corFour}

\begin{proof}
It is a simple consequence of proposition~2, corollary~3 and
the estimate~(11). The only thing we must do is to consider each
case in these results separately.
\end{proof}

\begin{note}
Although it will not be used in what follows, corollary~4 also holds when $c = \pm 1$.
The proof is similar: starting from other expressions for $K_{n}(x,\pm 1)$,
analogous results to proposition~2 and corollary~3 can be obtained,
and then corollary~4 follows.
\end{note}

We are now ready to extend theorem~1 to the general case $\gamma_{i} > -1$:

\begin{theoremFive}
Let $1 < p < \infty$.
If the inequalities
$$
\left\{
\begin{aligned}
A + (\alpha+1) \left(\frac{1}{p} - \frac{1}{2}\right) &< \frac{1}{4}, \\
B + (\beta+1) \left(\frac{1}{p} - \frac{1}{2}\right) &< \frac{1}{4}, \\
G_{i} + (\gamma_{i}+1) \left(\frac{1}{p} - \frac{1}{2}\right) &< \frac{1}{2}
\quad \forall i,
\end{aligned}
\right.
\quad
\left\{
\begin{aligned}
A + (\alpha+1) \left(\frac{1}{p} - \frac{1}{2}\right) &< \frac{\alpha+1}{2}, \\
B + (\beta+1) \left(\frac{1}{p} - \frac{1}{2}\right) &< \frac{\beta+1}{2}, \\
G_{i} + (\gamma_{i}+1) \left(\frac{1}{p} - \frac{1}{2}\right) &< \frac{\gamma_{i}+1}{2}
\quad \forall i,
\end{aligned}
\right.
$$
$$
\left\{
\begin{aligned}
a + (\alpha+1) \left(\frac{1}{p} - \frac{1}{2}\right) &> - \frac{1}{4}, \\
b + (\beta+1) \left(\frac{1}{p} - \frac{1}{2}\right) &> - \frac{1}{4}, \\
g_{i} + (\gamma_{i}+1) \left(\frac{1}{p} - \frac{1}{2}\right) &> - \frac{1}{2}
\quad \forall i,
\end{aligned}
\right.
\quad
\left\{
\begin{aligned}
a + (\alpha+1) \left(\frac{1}{p} - \frac{1}{2}\right) &> - \frac{\alpha+1}{2}, \\
b + (\beta+1) \left(\frac{1}{p} - \frac{1}{2}\right) &> - \frac{\beta+1}{2}, \\
g_{i} + (\gamma_{i}+1) \left(\frac{1}{p} - \frac{1}{2}\right) &> - \frac{\gamma_{i}+1}{2}
\quad \forall i,
\end{aligned}
\right.
$$
$$
A \leq a, \qquad B \leq b, \qquad G_{i} \leq g_{i}
\quad \forall i,
$$
hold, then there exists $C > 0$ such that
$$
\| S_{n}f \|_{L^{p}(u^{p}\,d\mu)}
\leq C \| f \|_{L^{p}(v^{p}\,d\mu)}
\quad \forall f \in L^{p}(v^{p}\,d\mu),
\quad \forall n \in \N.
$$
\end{theoremFive}

\begin{proof}
By induction on the number of negative exponents $\gamma_{i}$.
If $\gamma_{i} \geq 0$ $\forall i$, the result is true, as we saw
before (theorem~1).

Suppose there exist $k$ negative exponents $\gamma_{i}$,
with $k > 0$, and the result is true for $k-1$. Let $c \in (-1,1)$
be a point with a negative exponent~$\gamma$.
Let us remember the formula~(7):
$$
K_{n}(x,y) = (x-c) (y-c) K_{n-1}^{c}(x,y)
+ \frac{K_{n}(x,c) K_{n}(c,y)}{K_{n}(c,c)}.
$$
We define the operators
\begin{align*}
T_{n}f(x) &= \int_{-1}^{1}
\frac{K_{n}(x,c) K_{n}(c,y)}{K_{n}(c,c)}\, f(y) w(y)\,dy,\\
R_{n}f(x) &= \int_{-1}^{1}
(x-c) (y-c) K_{n-1}^{c}(x,y) f(y) w(y)\,dy.
\end{align*}
Then, $S_{n} = T_{n} + R_{n}$. We are going to study
firstly the operators~$T_{n}$:
$$
T_{n}f(x) = \int{K_{n}(x,c)}{K_{n}(c,c)}
\int_{-1}^{1} K_{n}(c,y) f(y) w(y)\,dy,
$$
thus
\begin{align*}
\| T_{n}f \|_{L^{p}(u^{p}w)}
&\leq \frac{\int_{-1}^{1}  | K_{n}(c,y) | v(y)^{-1} | f(y) | v(y) w(y) \,dy}
{K_{n}(c,c)}\, \| K_{n}(x,c) \|_{L^{p}(u^{p}w)} \\
&\leq \frac{\| K_{n}(x,c) \|_{L^{p}(u^{p}w)}
\| K_{n}(x,c) v(x)^{-1} \|_{L^{q}(w)}}{K_{n}(c,c)}\,
\| f v \|_{L^{p}(w)} \\
&= \frac{\| K_{n}(x,c) \|_{L^{p}(u^{p}w)}
\| K_{n}(x,c) \|_{L^{q}(v^{-q}w)}}{K_{n}(c,c)}\,
\| f \|_{L^{p}(v^{p}w)}.
\end{align*}
From corollary~4 it follows now that $\| T_{n}f \|_{L^{p}(u^{p}\,d\mu)}
\leq C \| f \|_{L^{p}(v^{p}\,d\mu)}$
$\forall f \in L^{p}(v^{p}\,d\mu)$,
$\forall n \in \N$. So, we only have
to proof the same bound for the operators~$R_{n}$.

But, if we denote by $S_{n}^{c}$ the partial sums of the Fourier
series with respect to the measure $(x-c)^{2}w(x)\,dx$,
it turns out that
$$
R_{n}f(x) = (x-c) \int_{-1}^{1} (y-c) K_{n-1}^{c}(x,y) f(y) w(y)\,dy
= (x-c) S_{n-1}^{c} \left(\frac{f(y)}{y-c}, x \right),
$$
whence
\begin{align*}
&\| R_{n}f \|_{L^{p}(u^{p}w)} \leq C \| f \|_{L^{p}(v^{p}w)}
\quad\forall f \in L^{p}(v^{p}w) \\
&\qquad\iff
\big\| (x-c) S_{n-1}^{c} \big(\tfrac{f(y)}{y-c}, x \big) \big\|_{L^{p}(u^{p}w)}
\leq C \| f \|_{L^{p}(v^{p}w)}
\quad\forall f \in L^{p}(v^{p}w) \\
&\qquad\iff
\| (x-c) S_{n-1}^{c}g (x) \|_{L^{p}(u^{p}w)}
\leq C \| (x-c) g \|_{L^{p}(v^{p}w)}
\quad\forall g \in L^{p}(|x-c|^{p} v^{p}w) \\
&\qquad\iff
\| S_{n-1}^{c}g \|_{L^{p}(|x-c|^{p}u^{p}w)}
\leq C \| g \|_{L^{p}(|x-c|^{p}v^{p}w)}
\quad\forall g \in L^{p}(|x-c|^{p} v^{p}w) \\
&\qquad\iff
\| S_{n-1}^{c}g \|_{L^{p}({\tilde{u}}^{p}(x-c)^{2}w)}
\leq C \| g \|_{L^{p}({\tilde{v}}^{p}(x-c)^{2}w)}
\quad\forall g \in L^{p}({\tilde{v}}^{p}(x-c)^{2}w),
\end{align*}
where $\tilde{u}(x) = | x - c |^{1-2/p} u(x)$
and $\tilde{v}(x) = | x - c |^{1-2/p} v(x)$.

Therefore, we must prove the boundedness of the partial sums
$S_{n}^{c}$ with the weights $( \tilde{u}, \tilde{v} )$.
But the Fourier series we are considering now corresponds to
the Jacobi generalized weight $(x-c)^{2}w(x)$, which has only $k-1$
negative exponents $\gamma_{i}$, since on the point $c$ the
exponent is $\gamma + 2 > 1$.
By hypothesis, the theorem holds in this case and we only have
to see that the conditions in the statement hold for the weights
$(x-c)^{2}w(x)$, $| x - c |^{1-2/p} u(x)$
and $| x - c |^{1-2/p} v(x)$.

Except for the point $c$, these weights have the same exponents
as $w$, $u$ and $v$; thus, those conditions are the same and, therefore,
they are satisfied. At the point $c$, the exponents are, respectively,
$$
\gamma + 2,\qquad g + 1 - 2/p,\qquad G + 1 - 2/p.
$$
So, the inequalities we have to check are the following:
\begin{align*}
(G + 1 - 2/p) + ( \gamma + 2 + 1 ) (1/p - 1/2 ) &< 1/2,\\
(G + 1 - 2/p) + ( \gamma + 2 + 1 ) (1/p - 1/2 ) &< ( \gamma + 2 + 1 )/2,\\
(g + 1 - 2/p) + ( \gamma + 2 + 1 ) (1/p - 1/2 ) &> - 1/2,\\
(g + 1 - 2/p) + ( \gamma + 2 + 1 ) (1/p - 1/2 ) &> - ( \gamma + 2 + 1 )/2,\\
G + 1 - 2/p &\leq g + 1 - 2/p.
\end{align*}
It is easy to see that all of them are satisfied, from our hypothesis.
Consequently, we get
$$
\| S_{n-1}^{c}g \|_{L^{p}({\tilde{u}}^{p}(x-c)^{2}w)}
\leq C \| g \|_{L^{p}({\tilde{v}}^{p}(x-c)^{2}w)}
\quad\forall g \in L^{p}({\tilde{v}}^{p}(x-c)^{2}w),
\quad\forall n \in \N;
$$
thus
$$
\| R_{n}f \|_{L^{p}(u^{p}w)} \leq C \| f \|_{L^{p}(v^{p}w)}
\quad\forall f \in L^{p}(v^{p}w), \quad\forall n \in \N.
$$
And, finally,
$$
\| S_{n}f \|_{L^{p}(u^{p}\,d\mu)} \leq C \| f \|_{L^{p}(v^{p}\,d\mu)}
\quad\forall f \in L^{p}(v^{p}\,d\mu), \quad\forall n \in \N.
$$
That is, the result is true for $k$ negative exponents~$\gamma_{i}$.
By induction, it is true in general and the theorem is proved.
\end{proof}

\begin{note}
It can be shown that the converse is also valid, that
is, the boundedness of the partial sums $S_{n}$ implies the
five conditions of the theorem.
\end{note}

%%%%%%%%%%%%%%%%%%%%%%%%%%%%%%%%%%
\begin{thebibliography}{MNT}
%%%%%%%%%%%%%%%%%%%%%%%%%%%%%%%%%%

\frenchspacing

\bibitem[B]{B}
\textsc{V. M. Badkov},
Convergence in the mean and almost everywhere of Fourier series
in polynomials orthogonal on an interval,
\textit{Math USSR Sb.} \textbf{24} (1974), 223--256.

\bibitem[MNT]{MNT}
\textsc{A. M\'{a}t\'{e}, P. Nevai, V. Totik},
Extensions of Szeg\H{o}'s theory of orthogonal polynomials.~II,
\textit{Constr. Approx.} \textbf{3} (1987), 51--72.

\bibitem[N]{N}
\textsc{P. Nevai},
``Orthogonal Polynomials'',
Memoirs of the Amer. Math. Soc., vol.~18,
no.~213, Providence, R.I., U.S.A., 1979.

\bibitem[R]{R}
\textsc{E. A. Rahmanov},
On the asymptotics of the ratio of orthogonal polynomials,
\textit{Math. USSR Sb.} \textbf{32} (1977), 199--213.

\bibitem[V]{V}
\textsc{J. L. Varona},
``Convergencia en L$^{p}$ con pesos de la serie de Fourier
respecto de algunos sistemas ortogonales'', Ph.~D. Thesis,
Sem. Mat. Garc\'{\i}a de Galdeano, sec.~2, no.~22. Zaragoza, 1989.

\end{thebibliography}

\end{document}
